% ============================================================
% CAPITOLO 1 — La Rivoluzione 0-Code: Vederla in Azione
% ============================================================
\chapter{La Rivoluzione 0-Code: Vederla in Azione}

\section*{Cosa Imparerai}

Alla fine di questo capitolo saprai:
\begin{itemize}
  \item Come si presenta una sessione di lavoro in 0-code (dal vivo)
  \item Perché il contesto è tutto e come cambia i risultati
  \item Cosa costruirai con questo libro (anteprima dei 9 progetti)
\end{itemize}

\begin{notebox}
Se non hai ancora letto l'Introduzione, fallo ora. Spiega il framework ADLC, il confronto con l'SDLC tradizionale, e le tre competenze fondamentali (Context Engineering, Risk Design, Confidence Tagging) su cui si basa tutto il libro.
\end{notebox}

% ---------------------------------------------------------
\section{Una Sessione di Lavoro in 0-Code: Dal Vivo}

La teoria la conosci dall'Introduzione. Ora vediamo come funziona \textbf{concretamente} --- che aspetto ha una sessione di lavoro in 0-code, cosa appare sullo schermo, cosa fai tu e cosa fa l'IA.

\subsection{Il setup: 30 secondi}

Apri VS Code. Apri una cartella vuota. Apri la chat di Copilot in Agent Mode. Scrivi:

\begin{lstlisting}
Crea un'applicazione Python che converte temperature tra Celsius,
Fahrenheit e Kelvin. L'utente sceglie la conversione da un menu
interattivo nel terminale. Aggiungi un file di test con pytest.
\end{lstlisting}

\subsection{Cosa succede nei 60 secondi successivi}

Copilot in Agent Mode:
\begin{enumerate}
  \item \textbf{Crea} il file \texttt{converter.py} con le funzioni di conversione
  \item \textbf{Crea} il file \texttt{test\_converter.py} con i test unitari
  \item \textbf{Esegue} \texttt{pytest} nel terminale integrato
  \item \textbf{Legge} l'output: 6 test passati su 6
  \item \textbf{Ti notifica}: ``Ho creato i file e i test passano tutti. Vuoi che aggiunga qualcosa?''
\end{enumerate}

Tu hai scritto 3 righe di istruzioni. L'IA ha generato $\sim$80 righe di codice funzionante e testato.

Ma ecco il punto: il codice è \textbf{generico}. Funziona, ma non segue una struttura specifica, non ha gestione degli errori sofisticata, non ha documentazione interna.

\subsection{Lo stesso task, con contesto}

Ora immagina di scrivere prima un file \texttt{\_CONTEXT.md}:

\begin{lstlisting}[language=yaml]
# Progetto: Temperature Converter

## Scopo
CLI tool per conversione temperature con focus su precisione scientifica.

## Vincoli
- Kelvin non puo' essere negativo (0 K = zero assoluto)
- Arrotondamento sempre a 2 decimali
- Messaggi di errore in italiano
- Naming: snake_case per tutto

## Struttura
converter/
|-- main.py          <- Entry point con menu
|-- conversions.py   <- Funzioni pure di conversione
|-- validators.py    <- Validazione input
+-- tests/
    |-- test_conversions.py
    +-- test_validators.py
\end{lstlisting}

E poi scrivi a Copilot:

\begin{lstlisting}
Leggi il _CONTEXT.md e implementa l'applicazione completa
seguendo tutte le specifiche.
\end{lstlisting}

Il risultato è \textbf{radicalmente diverso}: codice separato in moduli, validazione dello zero assoluto, test specifici per i casi limite, messaggi in italiano, struttura pulita.

\textbf{Stesse 3 righe di richiesta. Risultati incomparabili.} La variabile che cambia è il contesto.

% ---------------------------------------------------------
\section{Il Contesto è Tutto}

Questo è il principio più importante dell'intero libro, e merita di essere cristallizzato con un esempio concreto della differenza.

\textbf{Richiesta vaga:}
\begin{lstlisting}
Crea un'API per gestire gli utenti
\end{lstlisting}
\textbf{Risultato}: codice generico, probabilmente senza validazione, senza autenticazione, con naming inconsistente.

\textbf{Richiesta con contesto:}
\begin{lstlisting}
Crea un'API REST con Express.js per la gestione utenti.
Struttura: src/routes/, src/controllers/, src/models/.
Database: PostgreSQL con Prisma ORM.
Endpoint: GET /users, GET /users/:id, POST /users,
  PUT /users/:id, DELETE /users/:id.
Validazione input con Zod.
Risposte JSON standard: { success: boolean, data: T,
  error?: string }.
Gestione errori centralizzata con middleware.
\end{lstlisting}
\textbf{Risultato}: codice strutturato, validato, con pattern consistenti e pronto per la produzione.

La differenza non è nell'IA. È nel contesto. Nel Capitolo~3 imparerai a scrivere file \texttt{\_CONTEXT.md} professionali; nel Capitolo~4 inizierai a costruire il tuo primo progetto reale.

\begin{warnbox}
\textbf{Una nota sulle aspettative.} Questo libro non richiede di saper programmare, ma richiede \textbf{curiosità e disponibilità ad apprendere concetti nuovi}. A partire dal Capitolo~6 incontrerai termini come API, middleware, database, autenticazione, token JWT. L'IA scriverà tutto il codice, ma il tuo ruolo di \textit{architetto di contesto} richiede di comprendere \textit{cosa} stai chiedendo di costruire --- non \textit{come} scriverlo. Nell'Introduzione trovi un ``Crash Course'' con i concetti fondamentali, e ogni capitolo avanzato include un ``Box Teoria'' con le spiegazioni necessarie. Se un termine ti è oscuro, consulta il Glossario (Appendice~A).
\end{warnbox}

% ---------------------------------------------------------
\section{Cosa Costruirai con Questo Libro}

Questo libro è strutturato come una \textbf{progressione di progetti reali}, ognuno più complesso del precedente.

\subsection*{Progetto 1 --- Hello World (Capitolo 4)}
Il tuo primo programma generato interamente dall'IA. Imparerai a formulare richieste, leggere l'output e iterare.

\subsection*{Progetto 2 --- CLI Tool (Capitolo 5)}
Un'applicazione da riga di comando completa: parsing CSV/JSON, argomenti, test automatici. Il tuo primo \texttt{\_CONTEXT.md} strutturato.

\subsection*{Progetto 3 --- REST API (Capitolo 6)}
Il tuo primo server web con Express.js: endpoint CRUD, validazione, documentazione Swagger.

\subsection*{Progetti 4--7 --- Applicazione Web Full-Stack con OAuth 2.0 (Capitoli 7--10)}
Il cuore del libro. Costruirai:
\begin{itemize}
  \item \textbf{Progetto 4} (Cap.~7): Backend Node.js/Express con PostgreSQL e~Prisma
  \item \textbf{Progetto 5} (Cap.~8): Autenticazione OAuth~2.0 (Google/GitHub) con~JWT
  \item \textbf{Progetto 6} (Cap.~9): Frontend React con rotte protette e primo SKILL.md
  \item \textbf{Progetto 7} (Cap.~10): Integrazione end-to-end e dashboard CRUD completa
\end{itemize}

\subsection*{Progetti 8--9 --- App Mobile Flutter (Capitoli 11--12)}
\begin{itemize}
  \item \textbf{Progetto 8} (Cap.~11): Setup Flutter, prima app mobile e~SKILL.md Flutter
  \item \textbf{Progetto 9} (Cap.~12): Login OAuth da mobile, CRUD sincronizzato, deep link
\end{itemize}

\subsection*{Progetto finale --- Deploy Completo (Capitoli 13--15)}
Tutto in produzione: backend su cloud, frontend su Vercel, database PostgreSQL gestito, app mobile sugli store, CI/CD automatizzato.

% ---------------------------------------------------------
\section{Requisiti per Seguire Questo Libro}

\subsection*{Cosa devi sapere}
\begin{itemize}
  \item Usare un computer (Windows, macOS o Linux)
  \item Concetti base di come funziona un'applicazione (anche solo come utente)
  \item L'inglese tecnico di base (i nomi dei comandi e delle tecnologie sono in inglese)
\end{itemize}

\subsection*{Cosa NON devi sapere}
\begin{itemize}
  \item Non devi saper programmare in Python, JavaScript, Dart o qualsiasi linguaggio
  \item Non devi aver mai usato React, Node.js, Flutter o PostgreSQL
  \item Non devi aver mai usato un terminale o la riga di comando
  \item Non devi aver mai usato Git (ma lo imparerai strada facendo)
\end{itemize}

\subsection*{Cosa ti serve}
\begin{itemize}
  \item Un computer con Windows 10/11, macOS 10.15+ o Linux
  \item Una connessione internet stabile
  \item Un account GitHub (gratuito)
  \item Un abbonamento a GitHub Copilot OPPURE accesso a Claude Code
  \item VS Code installato (gratuito)
\end{itemize}

\begin{costbox}
GitHub Copilot\index{GitHub Copilot} Individual costa 10\$/mese. Claude Code\index{Claude Code} funziona a consumo tramite API Anthropic~\cite{anthropic:claude-code}. Il costo per completare tutti i progetti del libro è paragonabile a quello di un corso online ($\sim$20--50\,\euro). Nel Capitolo~2 vedremo le opzioni gratuite disponibili e come minimizzare la spesa.
\end{costbox}

% ---------------------------------------------------------
\section{La Promessa di Questo Libro}

Alla fine di questa lettura avrai:

\begin{enumerate}
  \item \textbf{Costruito} un'applicazione web completa con backend, frontend e autenticazione OAuth~2.0
  \item \textbf{Costruito} un'app mobile Flutter connessa al tuo backend
  \item \textbf{Deployato} tutto in produzione --- accessibile da chiunque nel mondo
  \item \textbf{Senza aver scritto manualmente} una singola riga di codice
\end{enumerate}

E soprattutto avrai acquisito un \textbf{metodo} --- il Context Engineering e l'ADLC --- che ti permetterà di costruire qualsiasi altra applicazione con lo stesso approccio.

Non è magia. È ingegneria del contesto. E inizia nel prossimo capitolo, con l'installazione del tuo ambiente di sviluppo.

\bigskip
\begin{quote}
\itshape ``Nell'era dello sviluppo 0-code, le parole umane perdono il loro ruolo di veicolo comunicativo informale e assumono il peso e le conseguenze di direttive in un linguaggio di programmazione procedurale ad alto livello.''
\end{quote}
